\documentclass[12pt]{article}
\usepackage[utf8]{inputenc}
\usepackage[T1]{fontenc}
\usepackage{hyperref}
\title{Silt'e Language Resources for K-12 Curriculum Development}
\author{Compiled by Claude Sonnet 5.5 (claude-sonnet-5-5) for Andr\'as Kornai}
\date{}
\begin{document}
\maketitle

\section{Introduction}
The code xst is marked ``Retired'' in the task prompt. The prompt gives no family, and its ``Bookkeeping'' family label suggests a retired code. This report covers the Silt'e language (also Silti, Selti, East Gurage) as found by one web search. Whether the retired code xst and the language described below are exactly the same variety was not established.

\section{Primary Educational Resources}
\begin{itemize}
\item Addis Ababa University thesis (June 2010), ``Silt'e as a Medium of Instruction''. It discusses Silt'e for classroom use, scripts, resources and teacher training. The search did not show its findings. \url{https://etd.aau.edu.et/items/86e3de77-027c-4b4a-b4f6-a98a1fe4334a}
\end{itemize}
No Silt'e textbook or primer was found.

\section{Linguistic and Reference Materials}
\begin{itemize}
\item Wikipedia overview of the Silt'e language: location in the Silt'e Zone, about 880,000 speakers in the 2007 census, and dialects. \url{https://en.wikipedia.org/wiki/Silt%CA%BCe_language}
\item Bush\"orf, a document on Selti (\emph{Encyclopaedia Aethiopica}, vol.~4) hosted on Scribd. It was not examined. \url{https://www.scribd.com/doc/47724774/Bustorf-EAE4-Selti}
\item A Silt'e--Amharic--English dictionary with concise grammar by E.~A.~Gutt (Addis Ababa, 1997) is cited in a scraped PDF. The citation is unreliable and the dictionary itself was not located. \url{https://etd.aau.edu.et/bitstreams/969ece9a-43ec-4152-9e9c-ba2186b27091/download}
\end{itemize}

\section{Religious and Community Resources}
The search found nothing specific.

\section{Digital and Audio Resources}
The search found no audio or software resources. Rosetta Project's page for xst is only a language-identification entry. \url{https://rosettapanglossia.longnow.org/index.php?title=Xst}

\section{Curriculum Development Context}
Silt'e is used in a region of Ethiopia with a sizeable population. The AAU thesis is a lead for the status of Silt'e as a medium of instruction. Ethiopian Ministry of Education and regional education bureau catalogues should be checked for textbooks.
\end{document}